
\section{The Pulse Broadening Function}
\label{sec:pulse_broadening_function}

Radio waves propagating through an inhomogeneous plasma can
follow multiple paths between a source and an observer.
Because these paths have different geometric lengths, photons
emitted simultaneously can arrive at different times.
The resulting distribution of arrival-time delays is
described by the \textbf{pulse broadening function} (PBF).

If the intrinsic pulse profile is $I_{\rm int}(t)$,
the observed profile is approximately the convolution

\begin{equation}
\boxed{
I_{\rm obs}(t)
=
\int_{-\infty}^{\infty}
I_{\rm int}(t')\,\mathrm{PBF}(t-t')\,dt'.
}
\end{equation}

We will derive the PBF for an idealized thin scattering
screen with an isotropic Gaussian distribution of
scattering angles.

\subsection{Geometric Time Delays}

Consider a source at distance $D$ from the observer,
with a thin scattering screen located a distance $D_s$
from the observer.

As derived previously, a ray arriving at an angle
$\theta$ relative to the direct line of sight
has an additional geometric path length

\begin{equation}
\Delta L
\simeq
\frac{D_sD}{2(D-D_s)}\theta^2.
\end{equation}

The corresponding time delay is

\begin{equation}
t=\frac{\Delta L}{c}.
\end{equation}

Defining the geometric factor

\begin{equation}
A\equiv
\frac{D_sD}{2c(D-D_s)},
\end{equation}

we can write

\begin{equation}
\boxed{t=A\theta^2.}
\end{equation}

This is the essential relationship: the geometric
delay is proportional to the square of the angular
offset.

Since $\theta^2\geq 0$, geometric scattering delays
are always non-negative relative to the direct path.

\subsection{The Angular Distribution of Scattered Rays}

Assume that the angular deflections in two transverse
directions, $\theta_x$ and $\theta_y$, are independent
Gaussian random variables with zero mean and equal
variance $\sigma_\theta^2$:

\begin{equation}
p(\theta_x,\theta_y)
=
\frac{1}{2\pi\sigma_\theta^2}
\exp\left[
-\frac{\theta_x^2+\theta_y^2}
{2\sigma_\theta^2}
\right].
\end{equation}

The total angular offset is

\begin{equation}
\theta^2=\theta_x^2+\theta_y^2.
\end{equation}

Because the distribution is isotropic, it is convenient
to transform into polar coordinates:

\begin{equation}
d\theta_x\,d\theta_y
=
\theta\,d\theta\,d\varphi.
\end{equation}

Integrating over the azimuthal angle $\varphi$,

\begin{align}
p_\theta(\theta)\,d\theta
&=
\int_0^{2\pi}
p(\theta_x,\theta_y)\,
\theta\,d\varphi\,d\theta\\
&=
\frac{\theta}{\sigma_\theta^2}
\exp\left(
-\frac{\theta^2}{2\sigma_\theta^2}
\right)d\theta.
\end{align}

Thus the probability density for the magnitude
of the angular offset is

\begin{equation}
\boxed{
p_\theta(\theta)
=
\frac{\theta}{\sigma_\theta^2}
\exp\left(
-\frac{\theta^2}{2\sigma_\theta^2}
\right),
\qquad \theta\geq0.
}
\end{equation}

This is a Rayleigh distribution.

\subsection{Converting Angular Deflections into Time Delays}

We now transform the angular probability distribution
into a distribution of arrival times.

Recall that

\begin{equation}
t=A\theta^2.
\end{equation}

Solving for $\theta$,

\begin{equation}
\theta=\sqrt{\frac{t}{A}}.
\end{equation}

Differentiating,

\begin{equation}
\frac{d\theta}{dt}
=
\frac{1}{2A\theta}.
\end{equation}

Conservation of probability requires

\begin{equation}
p_t(t)\,dt=p_\theta(\theta)\,d\theta.
\end{equation}

Therefore,

\begin{equation}
p_t(t)
=
p_\theta(\theta)
\left|\frac{d\theta}{dt}\right|.
\end{equation}

Substituting the angular distribution,

\begin{align}
p_t(t)
&=
\frac{\theta}{\sigma_\theta^2}
\exp\left(
-\frac{\theta^2}{2\sigma_\theta^2}
\right)
\frac{1}{2A\theta}\\
&=
\frac{1}{2A\sigma_\theta^2}
\exp\left(
-\frac{t}{2A\sigma_\theta^2}
\right).
\end{align}

Define the characteristic scattering timescale

\begin{equation}
\boxed{
\tau_{\rm scatt}
\equiv 2A\sigma_\theta^2.
}
\end{equation}

Using our expression for $A$,

\begin{equation}
\boxed{
\tau_{\rm scatt}
=
\frac{D_sD}{c(D-D_s)}
\sigma_\theta^2.
}
\end{equation}

The probability distribution of arrival times becomes

\begin{equation}
\boxed{
p_t(t)
=
\frac{1}{\tau_{\rm scatt}}
\exp\left(-\frac{t}{\tau_{\rm scatt}}\right),
\qquad t\geq0.
}
\end{equation}

For negative delays, the probability is zero.
We express this using the Heaviside step function,

\begin{equation}
H(t)=
\begin{cases}
0,&t<0,\\
1,&t\geq0.
\end{cases}
\end{equation}

Our final pulse broadening function is therefore

\begin{equation}
\boxed{
\mathrm{PBF}(t)
=
\frac{1}{\tau_{\rm scatt}}
e^{-t/\tau_{\rm scatt}}H(t).
}
\end{equation}

The normalization can be verified directly:

\begin{equation}
\int_{-\infty}^{\infty}
\mathrm{PBF}(t)\,dt
=
\int_0^\infty
\frac{1}{\tau_{\rm scatt}}
e^{-t/\tau_{\rm scatt}}dt
=1.
\end{equation}

\subsection{Why Scattered Pulses Have Exponential Tails}

Suppose the source emits an infinitely narrow pulse,

\begin{equation}
I_{\rm int}(t)=F\,\delta(t),
\end{equation}

where $F$ is the pulse fluence.

Convolving this with the PBF gives

\begin{align}
I_{\rm obs}(t)
&=
F\int_{-\infty}^{\infty}
\delta(t')\,\mathrm{PBF}(t-t')\,dt'\\
&=F\,\mathrm{PBF}(t).
\end{align}

Thus,

\begin{equation}
\boxed{
I_{\rm obs}(t)
=
\frac{F}{\tau_{\rm scatt}}
e^{-t/\tau_{\rm scatt}}H(t).
}
\end{equation}

An intrinsically instantaneous pulse is observed as
a one-sided exponential tail.

For a pulse with finite intrinsic width, the observed
profile is the convolution of the intrinsic emission
and this exponential response.

For example, if the intrinsic pulse is Gaussian,

\begin{equation}
I_{\rm int}(t)
=
\frac{F}{\sqrt{2\pi}\sigma}
\exp\left[
-\frac{(t-t_0)^2}{2\sigma^2}
\right],
\end{equation}

then convolution with the exponential PBF gives

\begin{equation}
\boxed{
\begin{aligned}
I_{\rm obs}(t)
={}&
\frac{F}{2\tau_{\rm scatt}}
\exp\left[
\frac{\sigma^2}{2\tau_{\rm scatt}^2}
-\frac{t-t_0}{\tau_{\rm scatt}}
\right]\\
&\times
\operatorname{erfc}\left[
\frac{1}{\sqrt{2}}
\left(
\frac{\sigma}{\tau_{\rm scatt}}
-\frac{t-t_0}{\sigma}
\right)
\right].
\end{aligned}
}
\end{equation}

Here $\operatorname{erfc}$ is the complementary
error function.

This is the familiar exponentially modified Gaussian
pulse profile, commonly used to model temporally
scattered radio bursts.

\subsection{Frequency Dependence}

The characteristic scattering timescale depends on
the angular distribution of the scattered radiation.

Since

\begin{equation}
\tau_{\rm scatt}\propto\sigma_\theta^2,
\end{equation}

the frequency dependence of the PBF timescale follows
from the frequency dependence of the scattering angle.

For an ideal Kolmogorov turbulent plasma in the
strong-scattering regime,

\begin{equation}
\sigma_\theta\propto\nu^{-11/5},
\end{equation}

and therefore

\begin{equation}
\boxed{
\tau_{\rm scatt}\propto\nu^{-22/5}
=\nu^{-4.4}.
}
\end{equation}

In practice, scattering timescales are often
parameterized relative to a reference frequency
$\nu_0$:

\begin{equation}
\boxed{
\tau_{\rm scatt}(\nu)
=
\tau_0
\left(\frac{\nu}{\nu_0}\right)^\alpha.
}
\end{equation}

The scattering index $\alpha$ can then be measured
from observations across multiple frequency bands.

\subsection{Limitations of the Exponential PBF}

The one-sided exponential PBF derived here relies on
an isotropic Gaussian distribution of scattering angles
and an idealized thin-screen geometry.

A real turbulent plasma does not necessarily produce
an exactly Gaussian angular distribution. In particular,
an ideal Kolmogorov phase structure function is not
quadratic, so its angular scattering distribution and
corresponding PBF are generally not exactly exponential.

Anisotropic scattering, multiple scattering screens,
finite screen sizes, and distributed scattering media
can also produce different pulse broadening functions.

Consequently, an exponential PBF is a useful model,
but it is not a universal description of radio-wave
scattering.
